% 第 4 章。数字：checks/check_ch04.py；引文：data/snapshots/money_facts.csv（逐字核对见
% data/fetch/verify_web_quotes.py）；美联储资产负债表、美国货币存量与政策利率：fred_fed_weekly、fred_money_monthly、
% fred_policy_rates_daily；中国人民银行资产负债表：pbc_balance_sheet；秘鲁：bcrp_dollarization；图：figsrc/fig_ch04.py。
\chapter{货币、银行与央行的资产负债表}\label{ch:money}

\begin{chapterguide}
\guidequestion{钱由谁创造？央行、商业银行和财政部的资产负债表怎样联动？}
\begin{guidegoals}
\item 写出央行、商业银行、财政部与家庭的简化资产负债表，追踪一笔贷款、一笔税款和一次国债发行；
\item 区分基础货币与广义货币，比较“贷款创造存款”与教科书货币乘数两种叙述；
\item 解释在准备金充足体系（地板体系）与走廊体系下，央行怎样设定短期利率；
\item 说出货币的三种职能，并举例说明它们可以分离。
\end{guidegoals}
\guideprereq{\chref{ch:accounts}\enspace\nameref{ch:accounts}；\chref{ch:prices}\enspace\nameref{ch:prices}}
\end{chapterguide}

2008 年 8 月到 12 月，美国的基础货币从 8,476 亿美元增加到 16,664 亿美元，四个月里几乎翻了一倍（1.97 倍）；同一时期，广义货币 M2 只从 7.76 万亿美元增加到 8.24 万亿美元，增长 6.3\%（图~\ref{fig:money_mult}）\cite{fred2026m}。按许多教科书的说法，银行拿到中央银行新发的钱，会按一个大致稳定的倍数放出贷款、创造存款，M2 应当跟着翻倍。它没有。六年以后，英格兰银行在自己的季刊上说，这种说法是一种误解。钱到底是谁创造的？中央银行又靠什么管住它？

这一章回答导引框里的问题：钱由谁创造？央行、商业银行和财政部的资产负债表怎样联动？依次讨论：2008 年秋天的美联储与 2014 年英格兰银行的说明（\secref{sec:money-world}）；四个部门的资产负债表，以及一笔贷款、一笔税款、一次国债发行怎样在其间流动（\secref{sec:money-sheets}）；货币乘数与“贷款创造存款”两种叙述（\secref{sec:money-creation}）；中央银行怎样设定短期利率：走廊与地板（\secref{sec:money-rate}）；货币的三种职能，以及人们为什么接受一张本身几乎不值钱的纸（\secref{sec:money-functions}）。\secref{sec:money-test} 用四组数据检验这些说法：2008 年以后的乘数、2019 年 9 月的货币市场、秘鲁的美元化，以及地板体系的成本。记号：$M$ 是货币存量（本章指广义货币 M2），$i$ 是利率（\appref{app:symbols}）。另有几个只在本章使用的记号：基础货币 $H$，公众手中的通货 $C$（不是\chref{ch:accounts}里的消费），银行的准备金 $R$，存款 $D$；$C/M$ 叫通货比率，$R/D$ 叫准备金比率，$m=M/H$ 是货币乘数。本章引用的外文文献，中文都是本书的翻译。

\section{当时的世界}\label{sec:money-world}

这一节交代：为什么“钱是谁创造的”在 2008 年以后成了需要重新回答的问题。先看美联储怎样在 2008 年秋天把资产负债表扩大了一倍半，又为什么要开始给银行的准备金付利息；再看英格兰银行 2014 年的那篇文章说了什么。

\subsubsection{2008 年秋天：资产负债表与准备金的利息}
2008 年 9 月 3 日（星期三），美联储的总资产是 9,053 亿美元；10 月 1 日是 15,040 亿美元；12 月 31 日是 22,395 亿美元，不到四个月扩大到原来的 2.47 倍（图~\ref{fig:money_fed}(a)）。扩大的部分不是国债：美联储持有的国债在这段时间里从 4,797 亿美元略降到 4,759 亿美元；增加的几乎全部是国债与抵押贷款支持证券以外的资产，从 4,255 亿美元增加到 17,635 亿美元。负债一边，银行存在美联储的准备金余额从 109 亿美元增加到 8,600 亿美元\cite{fred2026m}。联邦公开市场委员会（FOMC，美联储制定货币政策的委员会）10 月会议的纪要说，“现有流动性便利的扩大以及新便利的设立，使美联储的资产负债表显著扩大”，其中定期拍卖便利（TAF）的贷款余额在两次会议之间“增加了一倍多”\cite{fomc2008min}。

\fig{money_fed}{美联储的资产负债表，2007—2026 年，每周三，万亿美元。(a) 看资产的构成：美国国债（蓝色）、抵押贷款支持证券（浅蓝）与其他资产（浅褐，2008 年秋天主要是各种流动性便利）；总资产 2022 年 4 月 13 日最高，为 8.97 万亿美元，2026 年 10 月 7 日为 6.75 万亿美元。(b) 看负债的构成：联邦储备券（浅蓝）、准备金余额（蓝色）、财政部一般账户（紫色）、逆回购（最浅的蓝）与其他负债和资本（浅灰）。数据：美联储，H.4.1，经 FRED（序列 WALCL、TREAST、WSHOMCB、WLFN、WLODLL、WDTGAL、WLRRAL、WLTLECL、WCTCL，2026 年 10 月 14 日下载）。}{fig:money_fed}

10 月 6 日（星期一），美联储理事会宣布，将“为存款机构的法定准备金与超额准备金余额支付利息”。法定准备金是法规要求银行按存款的一定比例持有的准备金，这个比例叫法定准备金率（中国叫存款准备金率）；超过要求的部分叫超额准备金。为超额准备金付息的理由有两条：它“使美联储在用它的贷款项目应对信贷市场状况时有更大的余地，同时又能把联邦基金利率维持在联邦公开市场委员会设定的目标附近”；它还“使美联储能够按需要扩大资产负债表，提供维护金融稳定所需的流动性”\cite{fed2008ior}。法律早有授权，只是还没到生效的日子：2006 年的《金融服务监管减负法》原本允许美联储从 2011 年 10 月 1 日起付息，“刚刚通过的 2008 年《紧急经济稳定法》把生效日期提前到 2008 年 10 月 1 日”。法定准备金的利率定为每个维持期（按平均余额考核法定准备金的一段时间）内联邦基金目标利率的平均值减 10 个基点；超额准备金的利率起初定为维持期内最低的目标利率减 75 个基点；新规定从 10 月 9 日（星期四）起生效\cite{fed2008ior}。

联邦基金利率是银行之间隔夜拆借准备金的利率，FOMC 为它设定目标：2008 年 10 月 8 日从 2\% 降到 1.5\%，10 月 29 日降到 1\%，12 月 16 日改为一个区间，0 到 0.25\%\cite{fred2026m}。10 月 9 日，超额准备金利率是 $1.5\%-0.75\%=0.75\%$。美联储的说明写道：“为超额准备金付息，应当有助于给联邦基金利率建立一个下限。”\cite{fed2008ior}可是 10 月 9 日到 12 月 15 日，联邦基金有效利率（EFFR，市场成交的加权平均利率）每天都低于目标，平均比当时的超额准备金利率还低 43 个基点，最多低 89 个基点；这段时间里超额准备金利率的算法改过两次，10 月 23 日起是目标利率减 35 个基点，11 月 6 日起等于目标利率\cite{fred2026m}。10 月 28—29 日的 FOMC 会议纪要说，有效利率“持续低于 FOMC 的目标”，“不少机构显然愿意以低于超额准备金利率的利率拆出资金。非正式的报告表明，这尤其是那些没有资格从美联储的余额上拿到利息的机构”；付息“可能随着时间推移给联邦基金利率建立一个更有效的下限”\cite{fomc2008min}。这个漏水的下限，是\secref{sec:money-rate}的起点。

\subsubsection{2014 年：英格兰银行的说明}
英格兰银行的三位经济学家 McLeay、Radia 与 Thomas 在 2014 年第一季度的季刊上发表了《现代经济中的货币创造》。文章的摘要说：“现实中的货币创造，与一些流行的误解不同：银行并不只是充当中介，把储户存进来的钱贷出去；它们也不把中央银行货币‘乘’成新的贷款与存款。”\cite{mcleay2014}文章接着说，“在现代经济中，大部分货币以银行存款的形式存在”；“每当一家银行发放一笔贷款，它就同时在借款人的银行账户里创造一笔相应的存款，从而创造了新的货币”；“在正常时期，中央银行并不设定流通中的货币数量，中央银行货币也不会被‘乘’成更多的贷款与存款”\cite{mcleay2014}。

这些话与 2008 年的美国是吻合的：中央银行的钱翻了一倍，存款并没有跟着翻倍。文章也谈到英格兰银行自己的量化宽松：“作为量化宽松的副产品，新的中央银行准备金被创造出来。但它们不是传导机制的重要部分。”\cite{mcleay2014}文章说的是“正常时期”。中央银行在什么意义上“不设定”货币的数量，又靠什么约束银行，是\secref{sec:money-creation}与\secref{sec:money-rate}的内容。先要把账摆出来。

\section{四张资产负债表}\label{sec:money-sheets}

这一节回答：钱记在谁的资产负债表上？一笔贷款、一笔税款和一次国债发行，怎样在中央银行、商业银行、财政部与家庭之间流动？

\subsubsection{钱是谁的负债}
今天的钱都是某个机构的负债。纸币是中央银行的负债。银行在中央银行开的账户里的余额叫\textbf{准备金}，也是中央银行的负债，只有在中央银行开户的机构才能持有：英格兰银行的文章说，“准备金只能在银行之间借贷，因为消费者没有在英格兰银行开立准备金账户的途径”\cite{mcleay2014}。你在银行的存款是银行的负债：“银行存款只是银行欠它的客户多少钱的记录。所以它是银行的负债，而不是一种可以贷出去的资产。”\cite{mcleay2014}政府也在中央银行开户，美国财政部的账户叫财政部一般账户（TGA）；它的余额是中央银行对财政部的负债，不算在任何一种货币里。

\begin{intuition}[一张人人都收的借条]
你向朋友借钱，写一张借条；这张借条很难拿去买东西，卖东西的人不认识你，也不知道你会不会还。银行发放贷款，做的是一次借条的交换：你给银行一张借条（贷款合同），银行给你一张它自己的借条，也就是记在你账户上的存款。两张借条一样大，可银行的那张人人都收，可以拿去付账。银行的借条又要靠中央银行的借条（准备金）在银行之间结清，所以中央银行的负债站在最上面。你还清贷款，两张借条一起注销，那笔存款也就没有了。
\end{intuition}

中央银行对公众与银行的这两种负债，公众手中的通货与银行的准备金，合起来叫\textbf{基础货币}，也叫高能货币；公众手中的通货加上银行存款，叫\textbf{广义货币}。卡甘（Phillip Cagan，1927—2012）1965 年为美国国民经济研究局写的货币存量研究这样起头：“发行出来的高能货币，要么由公众作为银行以外的通货持有，要么由银行作为准备金持有”\cite{cagan1965}，于是
\begin{keyeq}
H=C+R,\qquad M=C+D.\label{eq:money-base}
\end{keyeq}
$M$ 包括哪些存款，各国的口径不同；本章的 $M$ 指 M2。两种货币的区别，见下面的对照框。

\begin{compare}{基础货币与广义货币}{基础货币 $H$}{广义货币 $M$}
\cmprow{是谁的负债}{中央银行：现钞与银行的准备金}{银行（存款）与中央银行（公众手中的现钞）}
\cmprow{谁能持有}{现钞人人可以持有；准备金只有在中央银行开户的机构能持有}{家庭、企业与其他非银行机构}
\cmprow{怎样增加}{中央银行买入资产或向银行放款}{银行发放贷款、买入债券或外汇；财政部支出（表~\ref{tab:money-steps}）}
\cmprow{规模}{美国 2026 年 8 月 5.41 万亿美元；中国 2026 年 6 月末的储备货币 40.6 万亿元}{美国 M2 23.28 万亿美元，是基础货币的 4.30 倍；中国 M2 356.7 万亿元，是储备货币的 8.79 倍}
\cmpsum{准备金是银行之间、银行与中央银行之间结账用的钱；广义货币是家庭与企业之间付账用的钱；公众手中的现钞两边都算。广义货币的大部分是银行的负债，中国 2026 年 6 月末流通中的现金（M0）只有 14.7 万亿元，占 M2 的 4.1\%。}
\end{compare}

对照框里的数字：美国 2026 年 8 月的基础货币是 54,116 亿美元，M2（未季节调整）是 232,843 亿美元\cite{fred2026m}；中国人民银行 2026 年 6 月末的资产负债表上，储备货币（中国人民银行对基础货币的叫法）是 40.6 万亿元\cite{pbc2026bs}，同期 M2 余额为 356.7 万亿元，“流通中货币（M0）余额为 14.7 万亿元”\cite{pbc2026q2}。

表~\ref{tab:money-fed} 是美联储 2026 年 10 月 7 日（星期三）的资产负债表。资产一边，国债与抵押贷款支持证券占 95.8\%，贷款只有 100 亿美元。负债一边，联邦储备券（即美元现钞）2.43 万亿美元，准备金余额 3.02 万亿美元，两者合起来是 5.46 万亿美元，差不多就是基础货币（基础货币另外还包括财政部发行的硬币）；准备金占总资产的 44.8\%。财政部一般账户有 8,858 亿美元。负债一边有一个负数：“其他负债”是 $-2{,}253$ 亿美元，它的来历在\secref{sec:money-test}。

\begin{table}[htbp]\centering\small
\caption{美联储的资产负债表，2026 年 10 月 7 日（十亿美元）}\label{tab:money-fed}
\begin{tabularx}{\linewidth}{@{}Y r Y r@{}}
\toprule
资产 & & 负债与资本 & \\\midrule
美国国债 & 4,566.4 & 联邦储备券（现钞） & 2,434.3 \\
\quad 其中短期国库券 & 562.2 & 逆回购协议 & 333.7 \\
抵押贷款支持证券 & 1,898.1 & 存款机构的存款（准备金余额） & 3,022.1 \\
贷款 & 10.0 & 财政部一般账户 & 885.8 \\
其他资产 & 273.1 & 外国官方存款 & 9.5 \\
 & & 其他存款 & 239.8 \\
 & & 其他负债（含递延资产） & $-225.3$ \\\midrule
 & & 负债合计 & 6,699.8 \\
 & & 资本 & 47.7 \\\midrule
资产合计 & 6,747.6 & 负债与资本合计 & 6,747.6 \\
\bottomrule
\end{tabularx}
\datasource{Federal Reserve, H.4.1（2026 年 10 月 8 日发布），表 5 与表 6；“其他资产”是总资产减去所列三项；“其他负债”是递延可用现金项目（0.3）、对信贷工具的财政部出资（0）与“其他负债与应计股息”（$-225.6$）之和。分项四舍五入，合计不一定相等。}
\end{table}

\subsubsection{一笔贷款、一笔税款、一次国债发行}
资产负债表的好处是能追踪：任何一笔交易都必须在两个以上的部门、每个部门的两边同时记账（\chref{ch:accounts}的四式记账）。表~\ref{tab:money-steps} 是一个只有四个部门的经济，依次发生五笔交易（教学数字），表里记的是每笔交易引起的变化。

\begin{table}[htbp]\centering\small
\caption{五笔交易在四张资产负债表上的变化（教学数字）}\label{tab:money-steps}
\begin{tabularx}{\linewidth}{@{}l l Y Y Y Y Y Y@{}}
\toprule
部门 & 项目 & （1）贷款 & （2）纳税 & （3）发债 & （4）支出 & （5）央行买债 & 合计 \\\midrule
家庭 & 存款（资产） & $+10$ & $-8$ & $-8$ & $+16$ & $+8$ & $+18$ \\
 & 国债（资产） & & & $+8$ & & $-8$ & 0 \\
 & 贷款（负债） & $+10$ & & & & & $+10$ \\
 & 净值 & & $-8$ & & $+16$ & & $+8$ \\\midrule
银行 & 准备金（资产） & & $-8$ & $-8$ & $+16$ & $+8$ & $+8$ \\
 & 贷款（资产） & $+10$ & & & & & $+10$ \\
 & 存款（负债） & $+10$ & $-8$ & $-8$ & $+16$ & $+8$ & $+18$ \\\midrule
中央银行 & 国债（资产） & & & & & $+8$ & $+8$ \\
 & 准备金（负债） & & $-8$ & $-8$ & $+16$ & $+8$ & $+8$ \\
 & 财政部存款（负债） & & $+8$ & $+8$ & $-16$ & & 0 \\\midrule
财政部 & 在央行的存款（资产） & & $+8$ & $+8$ & $-16$ & & 0 \\
 & 国债（负债） & & & $+8$ & & & $+8$ \\
 & 净值 & & $+8$ & & $-16$ & & $-8$ \\\midrule
\multicolumn{2}{@{}l}{广义货币 $\Delta M$} & $+10$ & $-8$ & $-8$ & $+16$ & $+8$ & $+18$ \\
\multicolumn{2}{@{}l}{基础货币 $\Delta H$} & 0 & $-8$ & $-8$ & $+16$ & $+8$ & $+8$ \\
\bottomrule
\end{tabularx}
\datasource{教学数字；计算见 \texttt{checks/check\_ch04.py}。为了简单，家庭不持有现钞，所以 $\Delta M$ 就是存款的变化，$\Delta H$ 就是准备金的变化。}
\end{table}

第 1 步，银行向家庭发放 10 的贷款。银行的资产（贷款）与负债（存款）同时增加 10，家庭的资产（存款）与负债（贷款）也同时增加 10。没有人事先存进一分钱，广义货币却多了 10；中央银行的账一笔没动。这就是英格兰银行说的“贷款创造存款”。

第 2 步，家庭纳税 8。家庭的存款减少 8；家庭的银行替它付款，把 8 的准备金转到财政部在中央银行的账户上，于是银行的存款与准备金各减少 8，中央银行的负债从准备金变成了财政部存款。纳税消灭了 8 的广义货币，也消灭了 8 的基础货币。财政部的净值增加 8，家庭的净值减少 8：税是一笔转移，没有对应的资产。

第 3 步，财政部发行 8 的国债，由家庭（例如一家养老基金）用存款买下。钱的流向与纳税一样，只是家庭换回了一张国债，净值不变。广义货币与准备金又各减少 8。

第 4 步，财政部支出 16，付给家庭（例如工资或福利）。财政部存款减少 16，银行的准备金与家庭的存款各增加 16。财政部收进来的钱，花出去时回到了银行体系。

第 5 步，中央银行从家庭手里买下那张国债。英格兰银行的文章用的就是这个例子：中央银行从一家养老基金买入 10 亿英镑的国债，“养老基金的银行在养老基金的账户上记入 10 亿英镑的存款，换取这些国债”，中央银行则在这家银行的准备金账户上记入 10 亿英镑\cite{mcleay2014}。在我们的表里，家庭的国债换成了存款，银行多了 8 的准备金。这就是量化宽松：广义货币与基础货币同时增加。

五笔交易合起来：广义货币增加 18，等于银行的贷款 10，加上财政赤字 8（支出 16 减税收 8），减去留在公众手里的国债（这里是 0）。基础货币增加 8，等于中央银行买进的国债。家庭的净值增加 8，恰好是财政部净值的减少，这就是\chref{ch:accounts}的部门余额恒等式：政府的赤字是私人部门的盈余。还要注意第三笔交易以后，银行的准备金累计减少了 16：如果它原来没有这么多准备金，就得向别的银行或中央银行借。银行在哪里、按什么价格借到准备金，是\secref{sec:money-rate}的问题。

\begin{example}{换一个买家}\label{ex:money-tax}
在表~\ref{tab:money-steps} 里，改由银行（而不是家庭）在第 3 步用准备金买下国债，并在第 5 步把它卖给中央银行，其余不变。（1）第 3 步与第 5 步的 $\Delta M$ 与 $\Delta H$ 各是多少？（2）五步合计，$\Delta M$、$\Delta H$ 与家庭的存款各是多少？与原表相比有什么不同？（3）回到原表，如果没有第 5 步，合计的 $\Delta M$ 与 $\Delta H$ 是多少？家庭手里多了什么？
\solution （1）第 3 步：银行的准备金减少 8，国债增加 8，存款不动，$\Delta M=0$，$\Delta H=-8$。第 5 步：银行的国债换成 8 的准备金，$\Delta M=0$，$\Delta H=+8$。银行卖债给中央银行，只是把一种资产换成另一种，不创造存款；非银行机构卖债，才同时创造存款。（2）$\Delta M=10-8+0+16+0=18$，$\Delta H=8$，家庭的存款增加 18，都与原表相同：国债最后都到了中央银行手里，中间经过谁，只影响中途的数字。（3）$\Delta M=10-8-8+16=10$，$\Delta H=0$；家庭多了 10 的存款与 8 的国债。这时财政赤字由家庭持有的国债来融资，量化宽松做的，就是把这 8 的国债换成存款与准备金。
\end{example}

\subsubsection{转账与结算}
家庭之间付款，钱也在资产负债表之间移动。甲在 A 银行有存款，付给在 B 银行开户的乙 5：A 银行对甲的负债减少 5，B 银行对乙的负债增加 5。存款换了主人，也换了银行。两家银行之间怎样了结？英格兰银行的文章说：“首先，买方的银行通过转移准备金与卖方的银行结算。”\cite{mcleay2014}在美国，这种转移经过美联储的两套系统。一套是逐笔、实时的：Fedwire 资金服务的参与者“从记入其联邦储备银行主账户的付款的最终性中得益”\cite{frbs2026}，钱一到账就不能撤回。另一套先轧差：全国结算服务（NSS）是一种“多边结算服务”，供那些替清算所、交易所等机构的参与者结算的存款机构使用，“借记与贷记是最终的、不可撤销的”，这是“以中央银行货币结算银行间债务”\cite{frbs2026}。Fedwire 的页面列出的用途里，还有“缴纳联邦税款”\cite{frbs2026}：表~\ref{tab:money-steps} 的第 2 步，在现实里可以经过这套系统完成。

\chref{ch:islamic}讲的汇票与轧差，是同一个道理：先把彼此的欠账抵消，只把净额用大家都信任的那种钱结清。今天那种钱就是中央银行的准备金。没有在中央银行开户的机构，通过它在另一家银行开的账户付款，后者叫\textbf{代理行}；跨国付款也是这样，靠各国银行在彼此那里开的账户层层转递。

\begin{deeper}[影响准备金的因素]
把表~\ref{tab:money-fed} 的两边移一下项，准备金就是中央银行的资产减去它的其他负债与资本：
\begin{keyeq}
R=A-N-D^{\mathrm T}-Q-X-K,\label{eq:money-reserves}
\end{keyeq}
其中 $R$ 是银行存在美联储的准备金余额，$A$ 是总资产，$N$ 是现钞（包括银行金库里的现钞），$D^{\mathrm T}$ 是财政部存款，$Q$ 是逆回购（2026 年 10 月 7 日的 3,337 亿美元里，3,314 亿美元来自外国官方与国际账户，货币市场基金等通过隔夜逆回购工具存进来的只有 23 亿美元），$X$ 是外国官方存款、其他存款与其他负债，$K$ 是美联储的资本（账面净值，与生产函数里的资本存量无关）。用 2026 年 10 月 7 日的数字（百万美元）：$6{,}747{,}560-2{,}434{,}312-885{,}783-333{,}739-23{,}944-47{,}716=3{,}022{,}066$，正是表里的准备金余额\cite{h41_2026}。美联储公布这张表的统计发布就叫“影响准备金余额的因素”（H.4.1）。式中由美联储自己决定的主要是 $A$（其中的贷款还要看银行来不来借）；公众要多少现钞、财政部账户里留多少钱、外国央行与货币市场基金把多少钱以逆回购的形式存在美联储，都由别人决定。它们任何一项增加，准备金就一对一地减少，除非美联储增加资产来抵消。财政部一般账户相当于准备金的 29.3\%：一个纳税日把 1,000 亿美元的税款转进这个账户，准备金就少 1,000 亿美元。2019 年 9 月的故事（\secref{sec:money-test}）就是从这里开始的。
\end{deeper}

\section{钱是怎样创造出来的}\label{sec:money-creation}

这一节回答：是中央银行先发出基础货币、银行再把它成倍地放大，还是银行先发放贷款、中央银行再供给所需的准备金？两种叙述用的是同一组恒等式，争的是因果的方向。

\subsubsection{货币乘数}
从式~\eqref{eq:money-base} 出发，把通货与准备金都写成比率：$C=(C/M)\,M$，$R=(R/D)\,D=(R/D)(1-C/M)\,M$。代入 $H=C+R$，整理得到
\begin{keyeq}
M=m\,H,\qquad m=\frac{1}{C/M+R/D-(C/M)(R/D)}.\label{eq:money-mult}
\end{keyeq}
这正是卡甘书中的式 (1)\cite{cagan1965}。他写道，货币的供给量“与通货比率和准备金比率反向变动，与高能货币的数量同向变动”\cite{cagan1965}。式~\eqref{eq:money-mult} 是恒等式，任何时候都成立。教科书的\textbf{货币乘数}叙述在它上面加了一个行为假设：两个比率大致稳定，于是中央银行决定 $H$，$M$ 随之按 $m$ 倍变动。英格兰银行的文章这样概括这种叙述：中央银行增加准备金，由于假定广义货币与基础货币之比是常数，这些准备金就被“乘”成大得多的银行存款与贷款的变化\cite{mcleay2014}。

卡甘本人并没有把这三个量当作原因。他说，它们“只是直接的决定因素，作为对货币存量的一种有用的分解，方便分析背后起作用的经济因素”；“这三个决定因素本身的表现并不能‘解释’货币供给的变化”\cite{cagan1965}。他还提醒，银行“不能同时控制准备金与负债的数量，只能控制两者的比率”；“银行可以同时增加贷款与存款，但一般不能同时增加货币准备金与存款”\cite{cagan1965}。

\begin{example}{2008 年的乘数}\label{ex:money-mult}
美国 2008 年 8 月的基础货币（月平均）为 8,476 亿美元，存款机构的总准备金为 458 亿美元，M2（未季节调整）为 77,582 亿美元；12 月三者分别为 16,664 亿、8,209 亿与 82,441 亿美元\cite{fred2026m}。把基础货币减去总准备金，近似当作公众手中的通货 $C$。（1）分别计算两个月的 $C$、$D$、$C/M$、$R/D$ 与乘数 $m$。（2）如果 12 月的乘数仍是 8 月的水平，M2 应当是多少？（3）哪个比率变了？这说明了什么？
\solution （1）8 月：$C=8{,}476-458=8{,}018$ 亿美元，$D=77{,}582-8{,}018=69{,}564$ 亿美元，$C/M=0.1033$，$R/D=0.00658$，$m=1/(0.1033+0.00658-0.1033\times0.00658)=9.15$，也就是 $77{,}582/8{,}476$。12 月：$C=8{,}455$ 亿美元，$D=73{,}986$ 亿美元，$C/M=0.1026$，$R/D=0.1110$，$m=4.95$。（2）$9.15\times16{,}664=152{,}530$ 亿美元，是实际 M2 的 1.85 倍。（3）通货比率几乎没变，准备金比率涨到原来的 16.9 倍：新增的基础货币几乎全部变成了银行存在美联储的准备金，没有变成贷款与存款。式~\eqref{eq:money-mult} 照样成立，只是 $m$ 从 9.15 降到 4.95，“乘数不变”的假设落空了。按式~\eqref{eq:money-mult} 的写法，也可以说银行“选择”了更高的准备金比率。可是准备金的总量是美联储买进资产与放款的结果，一家银行付出去的准备金总会落到另一家银行手里，整个银行体系只能像卡甘说的那样控制比率：要把比率压回去，就得多放贷款、多创造存款。10 月 9 日以后准备金有了利息，持有它的代价变了（\secref{sec:money-world}），银行没有理由急着这样做。
\end{example}

\fig{money_mult}{美国的基础货币、广义货币与两者之比，1959—2026 年。(a) 看 2008 年以后两条线怎样分开：M2（蓝色实线）与基础货币（褐色虚线），十亿美元，对数纵轴；浅蓝色竖条是 2008 年 8—12 月，基础货币增加到 1.97 倍，M2 只增长 6.3\%。(b) 看比率的跳动：M2 除以基础货币，1959—2007 年在 5.74 到 11.71 之间缓慢变动，2008 年 8 月是 9.15，12 月是 4.95，2014 年 8 月最低，为 2.81；2026 年 8 月是 4.30。月平均数，M2 未季节调整。数据：美联储，H.3 与 H.6，经 FRED（序列 BOGMBASE、M2NS，2026 年 10 月 14 日下载）。}{fig:money_mult}

\subsubsection{贷款创造存款}
英格兰银行的文章把因果倒了过来。“储蓄本身并不增加银行可以贷出去的存款或‘资金’”；“贷款的行为创造存款，这与教科书通常描述的顺序相反”\cite{mcleay2014}。准备金从哪里来？“在正常时期，准备金是英格兰银行按需供给商业银行的，换取它们资产负债表上的其他资产”；“中央银行今天通常不控制准备金的数量，而是设定准备金的价格，也就是利率”\cite{mcleay2014}。乘数理论要成立，“准备金的数量必须是贷款的一个有约束力的限制，中央银行必须直接决定准备金的数量”\cite{mcleay2014}，文章认为现实中这两条都不成立。于是，“对基础货币的这种需求，更可能是银行发放贷款、创造广义货币的结果，而不是原因”；“所以在现实中，货币乘数理论是反过来运作的”\cite{mcleay2014}。

那么什么限制银行创造货币？文章说有“三组主要的约束”\cite{mcleay2014}。第一是银行自己：“争夺贷款与存款的竞争，以及赚取利润的愿望，限制了银行创造货币。”第二是家庭与企业：它们可以用存款还债，“正像发放一笔新贷款创造货币，偿还银行贷款就消灭货币”。第三，“对货币创造的最终约束是货币政策”\cite{mcleay2014}：中央银行通过利率影响人们借多少钱，这一条留给\secref{sec:money-rate}。

\begin{compare}{货币乘数叙述与内生货币叙述}{货币乘数}{贷款创造存款（内生货币）}
\cmprow{因果的方向}{中央银行先供给基础货币，银行再成倍地放出贷款与存款}{银行先发放贷款、创造存款，再按需要取得准备金}
\cmprow{中央银行控制什么}{准备金的数量}{准备金的价格（短期利率）}
\cmprow{两个比率}{大致稳定，是行为常数}{由银行与公众的选择决定，随利率、监管与风险而变}
\cmprow{什么限制货币}{准备金与法定准备金率}{利润与竞争、借款人的需求、利率与监管}
\cmprow{2008 年的美国}{基础货币翻倍，M2 应翻倍；没有发生}{准备金多了，银行并不因此多放贷；吻合}
\cmpsum{两种叙述都满足式~\eqref{eq:money-mult}；分歧在于哪个量是中央银行定的、哪个量是随之调整的。中央银行盯住利率时，准备金的数量就是结果而不是原因。}
\end{compare}

\begin{pitfall}[银行把储户的存款贷出去]
这是对银行最常见的描述，支持它的最强论证来自单个银行的账本。一家银行发放贷款，借款人很快会把钱付给在别家银行开户的人，这家银行就要向对方转移准备金（\secref{sec:money-sheets}）。英格兰银行的文章自己也说，“如果一家银行全部用这种办法为新贷款融资，它很快就会耗尽准备金”；“通过吸收新的存款，银行可以在不减少准备金的情况下增加贷款”\cite{mcleay2014}。所以对一家银行来说，存款确实是放贷的资金来源，银行之间确实在争夺存款。反对它的最强论证来自整个银行体系：一家银行流失的准备金流进了另一家银行，存款的总量并不因为放贷而减少，反而因为放贷而增加（表~\ref{tab:money-steps} 的第 1 步）；存款是银行的负债，“不是一种可以贷出去的资产”\cite{mcleay2014}。两种说法讲的是不同层面的事：单个银行要为贷款找资金，银行体系作为整体却是在放贷时创造存款。把单个银行的经验推到整个体系，是一种\textbf{合成谬误}（fallacy of composition）：对每个个体成立的事，对整体未必成立。\chref{ch:accounts}的例题“节俭的一年”是同一类问题：家庭想多储蓄，整个经济的结果却未必如此。
\end{pitfall}

\subsubsection{争议：准备金要求还起不起作用}
两种叙述在法定准备金上分歧最明显。乘数叙述的最强证据来自中国。2011 年上半年，中国人民银行“6 次上调存款类金融机构人民币存款准备金率各 0.5 个百分点，累计上调 3 个百分点”，利用它“主动性较强、能够深度冻结流动性的特点”\cite{pbc2011q4}；当时的季度报告还公布货币乘数：2011 年年末，基础货币 22.5 万亿元，“货币乘数为 3.79，比上年末低 0.13。金融机构超额准备金率为 2.3\%”\cite{pbc2011q4}。2025 年 5 月，存款准备金率下调 0.5 个百分点，“向市场提供长期流动性约 1 万亿元”\cite{pbc2025q2}。如果准备金要求不起作用，中央银行为什么要这样用它？卡甘也提醒过另一面：“认为只有商业银行创造存款（或造成存款数量的变化）的看法，可能是误导的”\cite{cagan1965}，公众与中央银行的选择同样进入式~\eqref{eq:money-mult}。

内生货币叙述的最强证据来自美联储。2020 年 3 月 15 日，美联储理事会说，“多年来，法定准备金通过创造一个稳定的准备金需求，在货币政策的实施中起着中心作用”，然后“鉴于已经转向准备金充足的体系”，把法定准备金率降到零，3 月 26 日起生效\cite{fed2020rr}。法定准备金的作用，按这个说法，是稳定准备金的\emph{需求}，以便控制利率，而不是控制货币的数量。英格兰银行的文章也说，“法定准备金在今天大多数发达经济体的货币政策框架里不是一个重要方面”\cite{mcleay2014}。中国人民银行 2026 年的报告则这样描述 M2：它“主要是金融机构的负债、实体部门的资产，由银行通过投放贷款、增持债券、向非银融出、购买外汇等多种资产扩张行为派生而来”；“自上世纪 90 年代末以来，我国就已明确不再进行贷款规模管理”；2019—2025 年，“贷款增量与 M2 增量之比约从 107\% 降至 63\%”\cite{pbc2026q2}。这几句话说的正是“资产扩张创造存款”。本章引用的 2025—2026 年的三份报告里，也已经找不到“货币乘数”这个词。

两种读法可以对照着看同一组中国的数字。按同样的算法，2026 年 6 月末 M2 是储备货币的 8.79 倍，比 2011 年年末的 3.79 倍高得多；其间存款准备金率多次下调，6 月末超额准备金率只有 1.3\%\cite{pbc2026q2}。乘数叙述说，准备金率降了，乘数就升了，式~\eqref{eq:money-mult} 正是这样预言的。内生货币叙述说，2011 年提高准备金率，冻结的是购买外汇时投放出去的准备金（\secref{sec:money-rate}的案例框），起作用的是它对银行资金成本与利率的影响；乘数只是事后的比值。数据本身分不出谁对：恒等式在两种因果下都成立。要分辨，得看中央银行在准备金紧张时是否按它的利率供给准备金，这就是下一节的内容。

\section{央行怎样设定利率}\label{sec:money-rate}

\chref{ch:prices}把短期利率当作给定的：中央银行宣布一个数，市场就照着交易。这一节回答：中央银行如果不控制货币的数量，凭什么能决定短期利率？答案在准备金市场上。

\subsubsection{准备金市场：走廊与地板}
银行需要准备金，是为了彼此结算（\secref{sec:money-sheets}），有的国家还为了满足法定准备金要求。准备金短缺的银行，向准备金多余的银行隔夜借入，美国的这个市场叫联邦基金市场。中央银行在市场两边各放一个常设的窗口：一个是借款窗口，银行随时可以用抵押品按利率 $i_L$ 借到准备金；一个是存款窗口，或者为准备金付息，利率为 $i_D$。有了这两个窗口，没有银行愿意在市场上以高于 $i_L$ 的利率借钱，也没有能拿到 $i_D$ 的银行愿意以低于 $i_D$ 的利率拆出。所以
\begin{keyeq}
i_D\;\le\; i=f(R)\;\le\; i_L,\qquad f(R)\approx i_D\ \ \text{当}\ R\ge\bar R,\label{eq:money-corridor}
\end{keyeq}
其中 $f(R)$ 是准备金的需求曲线：准备金越少，银行愿意为它付的利率越高；准备金多到超过某个水平 $\bar R$，需求曲线就贴着 $i_D$ 变平。两个窗口的利率围成\textbf{利率走廊}。在\textbf{走廊体系}里，准备金是稀缺的，中央银行每天估计银行的需求，用公开市场操作调节供给，把市场利率引到走廊中间的目标上。在\textbf{地板体系}里，中央银行把准备金供给到 $\bar R$ 以上，市场利率贴着走廊的下沿；要改变利率，就直接改变 $i_D$。准备金的数量在两种体系里都由中央银行的资产负债表决定（式~\eqref{eq:money-reserves}），可是在地板体系里，数量的小幅变动不影响利率。

\subsubsection{美联储：从漏水的地板到充足准备金}
2008 年以后，美联储的准备金远远超过了需求曲线变平的位置，它从此进入地板体系，用准备金利率当作地板。可是这块地板一开始就漏水（\secref{sec:money-world}）：联邦基金市场上有一部分出借方没有资格从美联储拿利息，它们愿意以低于准备金利率的价格拆出。2009—2014 年，有效利率平均比准备金利率低 12.0 个基点，2014 年平均低 16.1 个基点（图~\ref{fig:money_rates}(a)）\cite{fred2026m}。到 2018 年，FOMC 的会议纪要说，联邦基金市场上的出借“目前以联邦住房贷款银行为主”\cite{fomc2018nov}：这个市场已经不再主要是银行之间的拆借。

2014 年 9 月，FOMC 公布了货币政策正常化的原则：美联储打算“主要通过调整超额准备金的利率”把联邦基金利率移入目标区间；“使用隔夜逆回购协议工具以及其他需要的补充工具”帮助控制利率；但隔夜逆回购工具“只在必要的程度上”使用，“在不再需要时逐步退出”；从长期看，美联储持有的证券“不多于有效实施货币政策所必需”\cite{fomc2014norm}。隔夜逆回购给没有准备金账户的货币市场基金等机构也开了一个存款窗口：它们可以按逆回购利率把钱借给美联储，就不会以更低的利率借给别人，于是逆回购利率成了第二块、更低的地板。2019 年 1 月，FOMC 说它打算继续在这样一个体系里实施货币政策：“充足的准备金供给保证了对联邦基金利率与其他短期利率水平的控制主要通过设定美联储的管理利率来实现，而且不需要对准备金供给的主动管理”\cite{fomc2019ample}。2020 年 3 月，法定准备金率降到零（\secref{sec:money-creation}）。

此后美联储的资产负债表经历了又一次扩张与收缩：总资产在 2022 年 4 月 13 日达到 8.97 万亿美元的峰值，准备金在 2021 年 12 月 15 日达到 4.27 万亿美元，逆回购在 2023 年 4 月 12 日达到 2.68 万亿美元（图~\ref{fig:money_fed}）\cite{fred2026m}。2025 年 10 月，FOMC 决定“在 12 月 1 日结束其证券总持有量的缩减”\cite{fomc2025oct}；12 月，它判断“准备金余额已经降到充足的水平”，将“根据需要开始购买较短期的国债，以持续维持充足的准备金供给”，同时把目标区间降到 3.5\%—3.75\%，准备金余额利率（IORB，2021 年 7 月起取代超额准备金利率）降到 3.65\%\cite{fomc2025dec}。2026 年 9 月 16 日，FOMC 把目标区间提高 1/4 个百分点，到 3.75\%—4\%；理事会把准备金余额利率提高到 3.90\%，9 月 17 日起生效；FOMC 指示纽约联储以 3.75\% 的利率开展隔夜逆回购，每个交易对手每天的上限为 1,600 亿美元，以 4.0\% 的利率开展常备的隔夜回购操作；一级信贷（贴现窗口）利率也提高到 4.0\%\cite{fomc2026sep}。逆回购的“逐步退出”，十二年以后还没有发生。

把这几个数排在一起，就是今天美联储几个管理利率的排列：逆回购 3.75\%（目标区间的下沿）$\le$ 有效利率 $\le$ 准备金余额利率 3.90\% $<$ 常备回购 4.0\%（目标区间的上沿）。2026 年 9 月 17 日到 10 月 9 日的 17 个营业日，有效利率每天都是 3.88\%\cite{fred2026m}。准备金余额利率与逆回购利率之差是 15 个基点：有效利率落在两块地板之间。2022—2024 年，逆回购的余额很大，有效利率平均只比逆回购利率高 3.0 个基点；2026 年逆回购余额缩小以后，这个差距是 13.3 个基点\cite{fred2026m}。在汇率浮动的经济里，中央银行日常盯住的就是这样一个短期利率；固定汇率下，它盯住的是汇率（\chref{ch:emcrises}）。\chref{ch:songming}拿南宋维持会子币值的办法与这两种目标作过比较。

\fig{money_rates}{美国的政策利率与准备金。(a) 看市场利率怎样被管理利率框住：联邦基金有效利率（蓝色实线），超额准备金利率与准备金余额利率（褐色虚线，2008 年 10 月 9 日起），隔夜逆回购利率（浅蓝色点线），FOMC 的目标（灰线，2008 年 12 月 15 日以前）与目标区间（灰带）；百分比。(b) 看准备金的需求曲线：横轴是每周三的准备金余额（万亿美元），纵轴是截至当天一周的有效利率减准备金利率的平均（基点），2009—2026 年；颜色由浅到深依次为 2009—2015、2016—2021、2022—2026 年；方块是 2019 年 9 月 18 日。数据：联邦基金有效利率与隔夜逆回购利率来自纽约联储，目标与准备金利率来自美联储理事会，准备金余额来自 H.4.1，均经 FRED（2026 年 10 月 14 日下载）。The Effective Federal Funds Rate is subject to the Terms of Use posted at newyorkfed.org. The New York Fed is not responsible for publication of the Effective Federal Funds Rate by this project, does not sanction or endorse any particular republication, and has no liability for your use.}{fig:money_rates}

\begin{example}{地板上的套利}\label{ex:money-floor}
2026 年 10 月 7 日，有效利率为 3.88\%，准备金余额利率为 3.90\%，逆回购利率为 3.75\%。货币市场的利息按一年 360 天计。（1）一家银行在联邦基金市场借入 10 亿美元隔夜资金，存在美联储拿准备金利息，每天赚多少？一年呢？（2）2014 年，有效利率平均比超额准备金利率低 16.1 个基点。同样的 10 亿美元，一年赚多少？（3）一家不能在美联储开准备金账户的货币市场基金，为什么不会以低于 3.75\% 的利率把钱借给别人？有效利率比它高多少？
\solution （1）利差 $3.90\%-3.88\%=2$ 个基点；每天 $10^9\times0.0002/360\approx556$ 美元，一年 20 万美元。（2）$10^9\times0.00161=161$ 万美元，约 160 万美元一年。（3）它可以按 3.75\% 把钱以逆回购的形式借给美联储，没有风险；所以逆回购利率是它的地板。有效利率比它高 13 个基点。套利为什么没有把利差消除？2008 年的会议纪要给出了一半答案：出借资金的机构没有资格拿准备金利息，只能把钱借给有资格的银行，后者压低了出价\cite{fomc2008min}。另一半在银行一边：这笔套利让银行的资产负债表膨胀 10 亿美元，而资产负债表并不是没有成本的；利差只有几个基点时，许多银行宁愿不做。所以地板体系里的市场利率贴着地板，却不正好落在地板上。
\end{example}

\begin{compare}{走廊体系与地板体系}{走廊体系}{地板体系}
\cmprow{准备金}{稀缺：中央银行每天估计需求、调节供给}{充足：供给超过需求曲线变平的位置 $\bar R$}
\cmprow{利率怎样定}{市场供求，中央银行把它引到走廊中间的目标上}{由管理利率（准备金利率）直接定，市场利率贴着地板}
\cmprow{资产负债表}{较小；扩张会压低利率}{较大；可以为别的目的扩张而不影响利率（2008 年的理由）}
\cmprow{例子}{中国人民银行：7 天期逆回购利率为锚，2026 年 6 月末超额准备金率 1.3\%}{美联储 2008 年以后；欧央行以存款便利利率引导（2024 年的框架重申）}
\cmprow{风险}{准备金需求的波动传到利率上}{准备金一旦低于 $\bar R$，利率突然上跳（2019 年 9 月）；要为大量准备金付息}
\cmpsum{两种体系用的是同一条需求曲线：走廊体系站在它倾斜的一段，地板体系站在它平坦的一段。}
\end{compare}

\subsubsection{欧央行与中国人民银行}
欧央行 2024 年 3 月 13 日宣布了新的操作框架：管理委员会“将继续通过存款便利利率（DFR）引导货币政策立场”，短期货币市场利率“预期在 DFR 附近变动，容忍一定的波动，只要它不模糊所要传递的政策信号”\cite{ecb2024}。主要再融资操作（MRO）的利率将调整到“与 DFR 之差从目前的 50 个基点缩小到 15 个基点”，2024 年 9 月 18 日起生效；边际贷款便利（MLF）与 MRO 之差“保持 25 个基点不变”；MRO 继续“以固定利率招标、全额分配”\cite{ecb2024}。理由是，“基于银行需要的、有弹性的中央银行准备金供给”最适合把流动性输送到整个欧元区的银行体系\cite{ecb2024}。这是一块地板（DFR），上面紧贴着一个按固定利率无限量供给准备金的窗口（DFR 加 15 个基点）。

中国人民银行的说法是，“调控货币市场短端利率是全球央行实施货币政策的普遍做法”，模式有两种：“一种是目标利率模式（以美联储为代表），央行直接以某个代表性市场利率（联邦基金有效利率 EFFR）作为政策利率”；“另一种是操作利率模式（以欧央行为代表）”\cite{pbc2026q2}。中国走的是后一种：报告“明确 7 天期逆回购操作利率是主要政策利率”；2024 年 7 月，7 天期逆回购“从竞争性利率招标调整为固定利率、充分满足一级交易商需求”；“MLF 利率完全退出政策利率功能”（这里的 MLF 是中期借贷便利，不是欧央行的边际贷款便利）；“2025 年起，短端利率调控目标从过去的 DR007 逐步转向 DR001”（DR001 是银行间隔夜回购利率）\cite{pbc2026q2}。2024 年 7 月增设的临时隔夜正、逆回购操作，利率定为政策利率减 20 个基点和加 50 个基点，“区间比超额准备金与常备借贷便利（SLF）的利率组成的利率走廊更窄”\cite{pbc2026q2}；2025 年全年，短期市场利率都运行在这个区间里\cite{pbc2025q4}。2026 年 6 月 17 日，这个区间改为 7 天期逆回购利率减、加各 25 个基点，“区间由此前的 70 个基点收窄为 50 个基点”\cite{pbc2026q2}。

政策利率本身，2025 年 5 月从 1.5\% 降到 1.4\%，全年只下调了这一次\cite{pbc2025q2,pbc2025q4}；2026 年上半年“保持不变”\cite{pbc2026q2}。按 1.4\% 计算，2026 年 6 月 17 日以后的操作区间是 1.15\%—1.65\%；外面一层是 SLF，“隔夜、7 天、1 个月 SLF 利率分别为 2.25\%、2.40\%、2.75\%”，隔夜 SLF 比政策利率高 85 个基点。6 月末，金融机构超额准备金率为 1.3\%；上半年 DR001 平均为 1.31\%，“略低于公开市场 7 天期逆回购操作利率”\cite{pbc2026q2}，低 9 个基点。超额准备金不多，市场利率贴着政策利率而不是贴着下沿：这是一个窄走廊，与美联储的地板正好相反。

\begin{case}{中国 · 2002—2026 · 从外汇占款到再贷款}
中国人民银行的资产负债表，资产一边经历了一次换血（图~\ref{fig:money_pbc}）。2002 年 12 月，总资产 5.11 万亿元，外汇占 43.3\%\cite{pbc2026bs}。此后国际收支长年顺差，中国人民银行买入外汇，每买一笔，就向卖外汇的银行投放一笔人民币准备金（表~\ref{tab:money-steps} 第 5 步的道理，只是买的是外汇而不是国债）。外汇占总资产的比重在 2014 年 3 月达到 83.3\%，外汇余额在 2014 年 5 月达到 27.30 万亿元\cite{pbc2026bs}。

投放出去的准备金太多，就要收回来。中国人民银行用了两种办法。一是发行中央银行票据，用一种不能当准备金用的负债换回准备金：资产负债表上“发行债券”一项在 2010 年 7 月达到 4.75 万亿元\cite{pbc2026bs}；2011 年“全年累计发行中央银行票据约 1.4 万亿元……截至 2011 年年末，中央银行票据余额约为 1.9 万亿元”\cite{pbc2011q4}。二是提高存款准备金率，把准备金锁在法定要求里：2011 年上半年连续 6 次上调（\secref{sec:money-creation}）。报告说，这些操作发生“在国际收支总体上继续保持较大顺差的背景下”\cite{pbc2011q4}；同一份报告也记下了转折的苗头：“第四季度外汇占款出现阶段性减少”\cite{pbc2011q4}。

2014 年以后，外汇不再增加，基础货币要从别处投放。资产负债表上“对其他存款性公司债权”（向银行的各种放款）占总资产的比重，从 2011 年 12 月的 3.6\% 升到 2016 年 12 月的 24.7\%、2026 年 8 月的 42.8\%（21.34 万亿元）；外汇降到 21.72 万亿元，占 43.6\%\cite{pbc2026bs}。中国人民银行说，再贷款“即中国人民银行向金融机构发放的贷款，是中国人民银行投放基础货币的一个渠道”\cite{pbc2025q4}。2026 年 8 月，总资产 49.86 万亿元，是 2002 年 12 月的 9.8 倍；储备货币 40.18 万亿元\cite{pbc2026bs}。

负债一边的基础货币，在这二十多年里一直增长；变的是它背后的资产。买外汇投放的钱，由国际收支决定，中国人民银行要用央票与准备金率去对冲；向银行放款投放的钱，由中国人民银行自己决定数量与利率，与上面说的利率走廊是一体的。外汇为什么会源源不断地流进来，又为什么在 2014 年以后流出，是\chref{ch:external}的问题。
\end{case}

\fig{money_pbc}{中国人民银行的资产负债表，2002—2026 年，月末，万亿元。(a) 看资产的构成怎样变：外汇（橙色）、黄金与其他国外资产（浅橙）、对政府债权（紫色）、对其他存款性公司债权（蓝色）与其他资产（浅灰）；外汇占总资产的比重在 2014 年 3 月达到 83.3\%。(b) 看负债的构成：货币发行（浅蓝）、其他存款性公司存款（蓝色，即银行的准备金）、发行债券（褐色，主要是中央银行票据，2010 年 7 月最高 4.75 万亿元；2017 年 6 月至 2018 年 10 月的表中这一项空白，按零画）、政府存款（紫色）与其他负债（浅灰）。资产与负债每月相等。数据：中国人民银行调查统计司，货币当局资产负债表（2026 年 10 月 14 日下载）。}{fig:money_pbc}

\section{货币的三种职能}\label{sec:money-functions}

这一节回答：什么东西算是钱？人们为什么愿意接受一张本身几乎不值钱的纸，或者银行账上的一个数字？

\subsubsection{三种职能可以分开}
货币通常被认为有三种职能：\textbf{交换媒介}，拿它买东西、付账；\textbf{计价单位}，用它标价、记账；\textbf{价值储藏}，用它把购买力留到将来。今天的人民币或美元同时承担这三种职能，所以人们很少想到它们可以分开。历史上它们常常是分开的。德川时代的日本，米是基本的记账单位，却不是交换媒介（\chref{ch:tokugawa}）。1990 年前后的恶性通胀以后，秘鲁的税仍按本币计算与缴纳，广义货币却有一半以上以外币计价（\secref{sec:money-test}）；恶性通胀中三种职能怎样分开，见\chref{ch:hyperinflation}。今天也有一种只在一个小圈子里充当交换媒介的钱：准备金只在银行与中央银行之间流通，家庭从来用不到它（\secref{sec:money-sheets}）。

既然职能可以分开，“一种东西为什么能当钱用”就要分开回答：它为什么被用来记账，为什么被接受为付款，为什么被人持有。对一张纸来说，最难回答的是后两个问题。两位相隔一百多年的作者给出了相近的回答。

\subsubsection{斯密：北美殖民地的纸币}
亚当·斯密（Adam Smith，1723—1790）在《国富论》（1776）第二篇第二章讨论了北美殖民地政府发行的纸币。这些纸币不付利息，要过若干年才兑付，殖民地政府却“宣布它是、并且在事实上使它成为按发行时的全部面值清偿债务的法定货币”\cite{smith1776}。斯密算了一笔账：在利率为 6\% 的国家，“15 年后支付的 100 英镑，只值比 40 英镑多一点的现钱”\cite{smith1776}。法律能强迫债权人接受它，却不能强迫卖东西的人：“成文法可以使一先令成为一基尼的法定货币，因为它可以命令法院解除提出这种偿付的债务人的债务；但没有一条成文法能迫使一个卖东西的人——他卖不卖完全自由——在价格上接受一先令等于一基尼。”\cite{smith1776}实际上，各殖民地纸币的价值相差悬殊：“100 英镑在有的殖民地有时被当作等于 130 镑当地货币，在另一些殖民地竟等于 1,100 镑当地货币”，差别来自“不同殖民地发行的纸币数量不同，以及最终兑付的期限的远近与可能性不同”\cite{smith1776}。

斯密接着指出纸币价值的另一个来源：税。“每个殖民地的纸币都可以按发行时的全部面值缴纳本省的税，它必然从这种用途中得到某种附加的价值，超过它仅凭最终兑付期限的实际或假想的远近所能具有的价值。”但附加价值有限：“这种附加价值的大小，取决于发行的纸币数量超过本殖民地缴税所能使用的数量的多少。在所有殖民地，它都远远超过了。”\cite{smith1776}他由此得出一个一般的命题：“一位君主如果规定，他的税收中有一定比例必须用某种纸币缴纳，就可能因此给这种纸币一定的价值，即使它最终兑付的期限完全取决于君主的意愿”；发行这种纸币的银行如果“小心地使它的数量总是略少于这样能够轻易使用的数量”，对它的需求甚至可能使它带有溢价\cite{smith1776}。

\begin{example}{斯密的 40 英镑}\label{ex:money-smith}
（1）按\chref{ch:prices}的贴现方法，利率为 6\% 时，15 年后支付的 100 英镑今天值多少？与斯密的“比 40 英镑多一点”相比如何？（2）一个当年就要缴 100 镑当地货币税款的殖民者，怎样看待手里一张面值 100 镑、15 年后才兑付的纸币？（3）斯密说，100 英镑在一些殖民地等于 130 镑当地货币，在另一些等于 1,100 镑。两种情形下，1 镑当地货币各值多少英镑？假定两地的纸币都是 15 年后按面值（当地的镑）兑付、利率都是 6\%，与（1）的贴现值相比，这说明了什么？（提示：当地的镑与英镑是同一个单位吗？）
\solution （1）$100/1.06^{15}=41.73$ 英镑，正是“比 40 英镑多一点”。（2）对他来说，这张纸今天就值 100 镑：它可以按面值抵 100 镑的税。只要他本来就要纳税，纸币的价值就不必等 15 年。（3）$100/130=0.77$ 英镑与 $100/1100=0.091$ 英镑，相差 8.5 倍。和 0.417 比较之前，先要看单位。斯密在同一段里说，发行最节制的宾夕法尼亚，纸币据说从没有跌到发行以前当地流通的金银以下；发行纸币以前，那里已经用法令让 5 先令英币当 6 先令 3 便士、后来当 6 先令 8 便士使用，所以即使流通的是金银铸币，100 英镑也合 133 镑当地货币，斯密说，当地的 1 镑即使是金银铸币，也比 1 英镑“低三成以上”，换成纸币以后，“很少比它低得多于三成”\cite{smith1776}。所以 130 这个比率很大一部分只是单位的差别：那里的纸币大致按面值与当地铸币等价流通，远高于“15 年后兑付”的贴现值（面值的 0.417），多出来的价值来自缴税等用途。在 1,100 的殖民地，同样面值的纸币只值前者的 $130/1100=0.12$；如果两地的铸币镑相差不大，它只值面值的一成多，连贴现值都不到：发行太多，缴税用不完，兑付本身也被怀疑了。税收给了纸币一个需求，发行的数量与人们对兑付的信心决定它最终值多少。
\end{example}

\subsubsection{克纳普：货币是法律的产物}
一百多年后，德国经济学家克纳普（Georg Friedrich Knapp，1842—1926）把这条线索推成了一种理论。他的《货币国定论》1905 年初版，1923 年出到第 4 版，1924 年有了据第 4 版的英文节译本\cite{knapp1924}。书的第一句话是：“货币是法律的产物。所以货币理论必须讨论法律史。”\cite{knapp1924}什么算是一个国家的货币？克纳普在 1895 年就提出，一国的货币“不是被强制普遍接受的东西，而是在公共收款处被接受的东西”；书中又说，“我们以这种货币在向国家的机构付款时被接受作为检验标准，就最贴近事实”；“决定性的不是发行，而是我们所说的接受”。他特意排除了“由国家发行”这个标准，“因为那会把极为重要的几种货币排除在外；我指的是银行券”\cite{knapp1924}。

今天的钞票本身几乎一文不值，却能买到东西：它是中央银行的负债，国家用它收税，法律规定它可以清偿债务；\chref{ch:songming}就从这里讲起。用克纳普的标准看表~\ref{tab:money-steps}：财政部收税时接受银行存款（第 2 步），这笔付款最后以中央银行的准备金完成（Fedwire 的“缴纳联邦税款”）；存款与准备金都在“公共收款处被接受”。可是两人都没有说，国家接受它缴税，就能决定它值多少：克纳普的标准回答的是“什么是货币”，斯密的附加价值则取决于发行的数量。税收与法律给一种货币一个需求的底；它值多少、人们愿意持有多少，还要看发行的数量与人们对它的信任。\secref{sec:money-test} 的秘鲁是一个例子。

\section{数据检验与反例}\label{sec:money-test}

这一章的几个说法都可以拿数据检验：货币乘数稳定吗？地板体系是不是只要设定利率就够了？国家规定用什么缴税，是不是就决定了人们持有什么货币？最后看地板体系的一笔成本。

\subsubsection{乘数稳定吗}
1959—2007 年，美国 M2 与基础货币之比在 5.74（1959 年 1 月）与 11.71（1985 年 2 月）之间缓慢变动（图~\ref{fig:money_mult}(b)）。它从来不是常数，只是变得慢，慢到可以近似当作一个稳定的比率。2008 年，它在四个月里从 9.15 降到 4.95（例~\ref{ex:money-mult}），2014 年 8 月降到最低的 2.81，这时准备金比率 $R/D$ 是 0.278；2026 年 8 月是 4.30，通货比率 0.106，准备金比率 0.141\cite{fred2026m}。2020 年 3 月以后法定准备金率是零，银行持有的准备金却仍相当于存款的 14\%：这不是法定要求。英格兰银行的文章列举了银行持有准备金的几个理由：应付公众提款、向别的银行付款、满足流动性监管的要求\cite{mcleay2014}；准备金有利息，又使持有它几乎没有代价。乘数作为恒等式一直成立；作为行为常数，它在 2008 年以后失效了。这支持内生货币的叙述，但也要看清它的范围：2008 年以后美联储处在地板体系，准备金有利息，持有准备金几乎没有代价；准备金稀缺、不付利息的年代，乘数变得慢得多，货币主义者依据的就是那时的经验（\chref{ch:monetarism}）。

\subsubsection{2019 年 9 月：地板需要足够的准备金}
美联储研究人员 2020 年的一篇说明回顾了这次事件。准备金“从 2007 年年底的不到 100 亿美元，增加到 2014 年 10 月约 2.8 万亿美元的峰值”，此后随着美联储缩减资产负债表而下降；2019 年 9 月中旬，“由于税款流出与国债的净发行”，准备金降到“不到 1.4 万亿美元”的多年低点\cite{anbil2020}。“2019 年 9 月中旬，在公司纳税日与国债净发行增加带来的准备金大幅下降中，隔夜货币市场利率跳升，并剧烈波动。”\cite{anbil2020}原因有两个：“首先，9 月 16 日到期的季度公司所得税从银行与货币市场基金的账户里被提走，进了财政部在美联储的账户”；“其次，9 月 16 日有 540 亿美元的长期国债完成结算”。“纳税与国债拍卖的结算抽走了大量现金，银行体系的准备金在两个营业日里减少了约 1,200 亿美元。”\cite{anbil2020}这正是式~\eqref{eq:money-reserves}：财政部存款增加，准备金一对一地减少。

利率的反应很剧烈。“9 月 16 日（星期一），SOFR 为 2.43\%，比前一个营业日高 13 个基点。回购市场的压力外溢到联邦基金市场，EFFR 为 2.25\%，比星期五高 11 个基点，处在 FOMC 目标区间的上沿”；“9 月 17 日，EFFR 升到目标区间上沿以上，为 2.3\%，SOFR 升到 5\% 以上”\cite{anbil2020}。SOFR 是有担保的隔夜融资利率，即以国债为抵押的隔夜回购利率。当时的超额准备金利率是 2.10\%，17 日的有效利率比它高 20 个基点；9 月 18 日（星期三）的准备金余额是 1.385 万亿美元，2019 年的最低值，截至当天一周的有效利率平均比准备金利率高 11.4 个基点\cite{fred2026m}。为什么准备金多的银行没有把钱借出去赚这笔高利？说明的作者认为，“不确定性的增加，以及利率上跳是暂时的，可能使银行不愿增加超额准备金的拆出”\cite{anbil2020}。

把各年放在一起，就是图~\ref{fig:money_rates}(b) 的准备金需求曲线。2014 年，准备金在 2.25 万亿到 2.82 万亿美元之间，有效利率平均比准备金利率低 16.1 个基点；2019 年，准备金在 1.39 万亿到 1.74 万亿美元之间，平均高 2.4 个基点；2026 年，准备金在 2.88 万亿到 3.18 万亿美元之间，平均低 1.7 个基点\cite{fred2026m}。准备金越少，利差越高，这与式~\eqref{eq:money-corridor} 的需求曲线一致；2019 年 9 月正是曲线变陡的地方。可是曲线本身会移动：2026 年各周的利差在 $-1$ 到 $-3$ 个基点之间，对应的准备金是 3 万亿美元左右；2018—2019 年，准备金在 1.75 万亿到 1.85 万亿美元之间的 11 周，利差在 $-3$ 到 0 个基点之间，平均 $-1.0$。同样的利差，2026 年要多出 1 万亿美元以上的准备金。2016—2026 年逐周的相关系数只有 $-0.29$。银行的规模、监管与付款的习惯都在变，$\bar R$ 也跟着变，而且只有在准备金接近它时才看得出来。2025 年 12 月 FOMC 判断准备金已降到“充足的水平”、开始购买短期国债，可以读作这样一种判断。地板体系不是“只设利率、不管数量”：数量要保持在一个事先并不确切知道的门槛以上。

\subsubsection{秘鲁：用本币纳税，却持有美元}
\chref{ch:silver}说，明代一条鞭法要求用银纳税，提高了人们必须持有的白银。一个更强的说法是：国家规定用什么缴税，就决定了人们持有什么货币。照这个说法，用本币纳税的国家就不该有大量的外币存款。秘鲁是一个反例。1990 年 8 月，秘鲁利马的消费者价格 12 个月里上涨了 12,377\%；1990 年 12 月的 12 个月涨幅是 7,650\%，1991 年 12 月是 139.2\%；1997 年 2 月起降到 10\% 以下，此后再没有超过，2001 年以后最高是 2022 年 6 月的 8.8\%（图~\ref{fig:money_peru}(a)）\cite{bcrp2026}。通胀过去了，美元却没有走：存款性公司的广义货币中以外币计价的比重，1992 年 1 月是 59.4\%，1998 年 12 月是 70.1\%，2000 年 5 月达到最高的 73.2\%，而那个月的 12 个月通胀率只有 3.3\%（图~\ref{fig:money_peru}(b)）\cite{bcrp2026}。

秘鲁的税并不收美元。2002 年 9 月 24 日的第 151-2002-EF 号最高法令在前言里引述现行的规定：账簿必须“以西班牙语记载、以本国货币表示”，只有收到或进行以外币计价的外国直接投资的纳税人可以例外\cite{peru2002}；即使是这些用美元记账的企业，“为申报与缴纳税款，申报中需要考虑的每一个组成部分都必须换算成本国货币”\cite{peru2002}。那个月，广义货币中外币的比重是 67.6\%\cite{bcrp2026}。

此后的下降是缓慢的：2010 年 12 月 37.6\%，2013 年 12 月 32.9\%。中央银行一直担心美元化的风险，国际清算银行的一篇工作论文说，在秘鲁这样“一个小型的、部分美元化的经济”里，“本币大幅贬值，会使收入与资产主要以本币计价的债务人要支付的本金与利息大大增加”\cite{minaya2017}。2013 年起，秘鲁中央银行实行去美元化计划，“通过以信贷的变化为条件、以外币缴存的额外准备金，抑制信贷美元化”：外币信贷的增长超过 10\%、15\% 与 20\% 三道限额的银行，分别要多缴 1.5\%、3\% 与 5\% 的准备金；到 2015 年 6 月，外币信贷总额（不含外贸业务）要降到 2013 年 9 月的 95\% 以下\cite{minaya2017}。论文发现，这一计划实施以后，“住房贷款的美元化显著下降”\cite{minaya2017}。计划针对的是贷款；同一时期，广义货币中外币的比重从 32.9\% 升到 2015 年 12 月的 37.4\%，可这个比重是把外币部分按月末汇率折成索尔计算的。这两年里索尔从 2.80 贬到 3.41 索尔兑 1 美元，同样多的美元折成索尔就多了：以美元计的外币部分与本币部分都增长了 14.0\%，按 2013 年 12 月的汇率折算，2015 年 12 月的比重仍是 32.9\%\cite{bcrp2026}。2025 年 12 月，它降到最低的 25.4\%；2026 年 8 月是 26.0\%\cite{bcrp2026}。

秘鲁的经验说明：用本币纳税确实给了本币一个需求，纳税的人总得在缴税时持有本币；可是持有多少，不由税收决定：恶性通胀结束以后的十年里（1992—2001 年），外币的比重一直在 55.1\% 到 73.2\% 之间，二十多年的低通胀与主动的政策才把它降到四分之一左右。这正是斯密的提醒：税收给了一种货币“附加的价值”，但不能保证它的价值。所以那个更强的说法站不住；更稳妥的说法是，税收给了一个需求的底，数量与信任决定其余。

\fig{money_peru}{秘鲁的通胀与美元化。(a) 看物价的那次爆发：利马消费者价格指数的 12 个月涨幅，纵轴按 1 加通胀率取对数，1990 年 8 月为 12,377\%；1997 年 2 月以后一直低于 10\%。(b) 看通胀过去以后美元留了多久：存款性公司的广义货币中以外币计价的比重（\%，外币部分按月末汇率折成索尔），1992 年 1 月为 59.4\%，2000 年 5 月最高，为 73.2\%，2026 年 8 月为 26.0\%；灰色点线是第 151-2002-EF 号最高法令（2002 年 9 月 24 日；它重申申报纳税用本币），浅蓝色竖条是去美元化计划的第一阶段（2013 年 9 月至 2015 年 6 月）。数据：秘鲁中央储备银行，BCRPData（序列 PN01273PM 与 PN00025MM，2026 年 10 月 14 日下载）。}{fig:money_peru}

\subsubsection{地板的代价}
地板体系要为大量准备金付息。2026 年 10 月 7 日的准备金余额是 3.02 万亿美元，按 3.90\% 的准备金余额利率，一年的利息约为 1,179 亿美元。美联储把扣除成本后的收益上缴财政部；它的资产大多是过去买入的固定利率债券，当它为准备金与逆回购支付的利息超过资产带来的收入时，就没有收益可以上缴。H.4.1 的表 6 里，“应上缴财政部的收益”一行是 $-2{,}335.87$ 亿美元，注释说：“负数表示累积的递延资产头寸，它发生在收益不足以支付运营成本、股息与维持盈余的时期。”\cite{h41_2026}这个数是美联储资本的 4.9 倍；表~\ref{tab:money-fed} 里那个负的“其他负债”，主要就是它。以后美联储有了收益，要先抵掉这笔递延资产，才能再向财政部上缴。不过这 1,179 亿美元不全是地板本身的代价：准备金的另一边是美联储持有的债券，它们也有利息收入；资产负债表小一些，这些债券的利息就由财政部直接付给私人持有者。递延资产来自用随政策利率变动的负债去持有过去买入的低息长期债券，也就是利率风险；资产如果是短期国库券，利息的收与付大致相抵。

两种体系的利弊，FOMC 在 2018 年 11 月的会议上摆过一次，都记在纪要里\cite{fomc2018nov}。支持准备金充足体系的一方说，根据近年的经验，这种体系“在各种市场条件下都能很好地控制短期货币市场利率”；准备金有限时，流动性监管或风险偏好的变化会使准备金的需求“更难预测”，利率就难以控制；而且即使要为缓解流动性紧张注入大量准备金，或者在利率降到下限时大规模购买资产，它也不会失去对利率的控制。2008 年为准备金付息，理由之一正是让美联储能够“按需要扩大资产负债表”\cite{fed2008ior}。另一方列出了它的缺点：“难以准确确定这种体系所需的准备金数量，需要维持相当大的准备金与证券持有量，以及为准备金付息带来的相当大的持续利息支出”；有人还说，回到准备金有限的体系，可以证明美联储能够完全退出为应对危机而采取的政策，“从而可能提高公众对这些政策的接受程度或它们的效果”\cite{fomc2018nov}。2019 年 1 月，FOMC 选择了前者：利率的控制“主要通过设定美联储的管理利率来实现，而且不需要对准备金供给的主动管理”\cite{fomc2019ample}。

2019 年 9 月说明，地板体系也需要一块天花板。9 月 17 日纽约联储宣布回购操作、直接向市场供给准备金以后，第二天两个市场的利率分布就回到了接近全年平均的水平\cite{anbil2020}。2021 年 7 月，FOMC 设立常备回购便利，作为“货币市场的后备”\cite{fomc2021srf}，也就是今天那 4.0\% 的上沿。两种体系怎样取舍，取决于怎样权衡利率的平稳、资产负债表的灵活、这笔利息与中央银行在金融市场里的分量；资产负债表为什么会扩张到今天的规模，见\chref{ch:zlb}与\chref{ch:gfc}。

\section{今天还剩什么}\label{sec:money-legacy}

从 2008 年的美联储到今天的中国人民银行，这一章的几条线索逐条看：
\begin{itemize}
\item \textbf{准备金与结算。}存款是银行的负债，银行之间最终用中央银行的准备金结算。稳定币与中央银行数字货币之争，争的是这个层级里谁能直接持有中央银行的负债（\chref{ch:digitalmoney}）。
\item \textbf{货币的等级。}中央银行的负债在最上面，银行存款在它下面，再往下是各种承诺按面值兑换存款的负债；平时它们一样好用，危机时层与层之间的兑换会中断（\chref{ch:hierarchy}、\chref{ch:gfc}）。
\item \textbf{乘数与货币主义。}式~\eqref{eq:money-mult} 是卡甘分解美国货币存量的工具；1930 年代美国的通货比率与准备金比率怎样变化，是\chref{ch:depressionus}讨论大萧条时的一个中心问题；控制货币数量的主张，见\chref{ch:monetarism}。
\item \textbf{地板与量化宽松。}2008 年以后的地板体系，是量化宽松能够在不失去利率控制的情况下进行的前提；零利率下限附近的政策，见\chref{ch:zlb}。
\item \textbf{财政部与中央银行。}税款、国债发行与财政支出都经过中央银行的账户（表~\ref{tab:money-steps}）。谁为谁融资、中央银行能不能拒绝，是\chref{ch:fiscaldominance}的问题。
\item \textbf{税收与货币。}从一条鞭法（\chref{ch:silver}）到宋元明的纸币（\chref{ch:songming}），国家用什么收税，一直影响着人们用什么钱；秘鲁说明这种影响有边界。现代货币理论把克纳普的说法推得更远，它的主张与争议见\chref{ch:minsky}。
\end{itemize}

这一章只看了一个国家内部的账。中国人民银行的资产在 2011—2014 年约有八成是外汇：外汇从哪里来，为什么要攒这么多，汇率又由什么决定，是\chref{ch:external}的问题。

\begin{takeaways}
\begin{itemize}
\item 今天的钱都是负债：现钞与准备金是中央银行的负债，合起来是基础货币 $H=C+R$；存款是银行的负债，广义货币 $M=C+D$。2026 年 8 月，美国的 M2 是基础货币的 4.30 倍；2026 年 6 月末，中国的 M2 是储备货币的 8.79 倍。
\item 银行发放贷款时创造存款，偿还贷款时消灭存款；纳税与非银行机构买国债使存款流进财政部的账户，财政支出与中央银行向非银行机构买债使存款流回来。表~\ref{tab:money-steps} 的五笔交易后，广义货币增加 18，等于贷款 10 加财政赤字 8。
\item 货币乘数 $m=1/[C/M+R/D-(C/M)(R/D)]$ 是恒等式；“乘数稳定、中央银行决定基础货币”是一个行为假设，在 2008 年以后的美国失效了：乘数四个月里从 9.15 降到 4.95。法定准备金在中国仍是主动使用的工具，在美国 2020 年降到零；两种叙述的分歧在因果的方向。
\item 中央银行设定的是准备金的价格。走廊体系里准备金稀缺，市场利率被引到走廊中间；地板体系里准备金充足，市场利率贴着准备金利率。2026 年 10 月，美联储的逆回购利率 3.75\%、有效利率 3.88\%、准备金余额利率 3.90\%、常备回购 4.0\%；中国人民银行的 7 天期逆回购利率 1.4\%（2026 年 6 月末），临时操作区间上下各 25 个基点。
\item 地板需要足够的准备金：2019 年 9 月准备金降到 1.4 万亿美元以下，一个纳税日就使有效利率突破目标区间的上沿、SOFR 超过 5\%。地板也有成本：美联储为 3 万亿美元的准备金付息，2026 年 10 月累积的递延资产是 2,336 亿美元。
\item 货币的三种职能可以分开。税收与法律给一种货币一个需求的底（斯密的“附加价值”，克纳普的“接受”），数量与信任决定其余：秘鲁用本币纳税，2000 年 5 月广义货币中仍有 73.2\% 是外币。
\end{itemize}
\end{takeaways}

\begin{selfcheck}
\begin{enumerate}
\item 家庭用存款偿还 4 的银行贷款。四张资产负债表怎样变？广义货币与基础货币各变多少？
\item 通货比率 $C/M=0.1$，准备金比率 $R/D=0.02$，乘数是多少？基础货币为 100 时 $M$ 是多少？如果准备金比率升到 0.2，乘数是多少？要维持同样的 $M$，基础货币要有多少？
\item 某个纳税日，企业把 1,000 亿美元税款付进财政部一般账户，其他条件不变。按式~\eqref{eq:money-reserves}，准备金怎样变？美联储要使准备金不变，可以怎么做？这件事在地板体系与走廊体系里对利率的影响有什么不同？
\item 某中央银行的存款便利利率为 2.0\%，贷款便利利率为 2.5\%，政策目标为 2.25\%。在走廊体系与地板体系里，隔夜市场利率大约各在哪里？如果它把三个利率都提高 0.25 个百分点，两种体系里市场利率各变成多少？需要改变准备金的数量吗？
\item 利率为 5\% 时，10 年后按面值兑付的 100 镑纸币今天值多少？如果当年就可以按面值缴税，对一个当年要缴 100 镑税的人，它值多少？
\item 下列各是货币的哪种职能，由什么东西承担？（a）德川日本按米的石数计算武士的俸禄；（b）1999 年的秘鲁人把储蓄存成美元存款；（c）2002 年以后，秘鲁用美元记账的企业仍用本币缴税；（d）两家美国银行用准备金结清当天的轧差。
\end{enumerate}
\end{selfcheck}

\begin{answers}
\item 家庭的存款与贷款各减少 4；银行的贷款（资产）与存款（负债）各减少 4；中央银行与财政部不变。广义货币减少 4，基础货币不变：偿还贷款消灭货币。
\item $m=1/(0.1+0.02-0.002)=8.47$，$M=847$。准备金比率为 0.2 时，$m=1/(0.1+0.2-0.02)=3.57$；要维持 $M=847$，基础货币要有 $847/3.57=237$。
\item 财政部存款 $D^{\mathrm T}$ 增加 1,000 亿美元，准备金减少 1,000 亿美元。美联储可以买入 1,000 亿美元的证券（或向银行放款），使 $A$ 增加同样的数。在地板体系里，只要准备金仍在 $\bar R$ 以上，利率几乎不动；接近 $\bar R$ 时（2019 年 9 月），利率会跳。在走廊体系里，准备金一减少利率就上升，所以中央银行要在当天用公开市场操作抵消。
\item 走廊体系：在目标 2.25\% 附近；地板体系：在 2.0\% 附近，通常略低于它（\secref{sec:money-rate}）。三个利率都提高 0.25 个百分点后，前者约 2.50\%，后者约 2.25\%。两种体系都不需要改变准备金的数量：走廊体系里，走廊整体上移，需求曲线上同一数量对应的利率也上移；地板体系里，地板本身上移。
\item $100/1.05^{10}=61.39$ 镑。对当年要缴 100 镑税的人，它值 100 镑。
\item （a）计价单位，米；（b）价值储藏，美元；（c）计价单位与向国家付款的手段，本币；（d）银行之间的交换媒介（结算），中央银行的准备金。
\end{answers}
